\chapter{Use in an $N$-body code}\label{ch:nbody}

This chapter gives the equations to be implemented in an $N$-body integrator for the tide raised by the star in the planet, at degree two. It does not depend on a particular code. The tide raised in the star uses the same equations (Chapter~\ref{ch:star}), and the implementation in Posidonius is the subject of Chapter~\ref{ch:posidonius}. Every formula of this chapter was checked numerically against an independent computation of the exact linear response of the planet, which does not use Kaula's expansion (Section~\ref{sec:validation}).

\section{Inputs, outputs and frame of the force}\label{sec:io}
\paragraph{Inputs,} all in inertial axes, at the time $t$ of the force evaluation:
\begin{itemize}
  \item $\vr=\vr_{\mathrm p}-\vr_\star$ and $\mathbf v=\mathbf v_{\mathrm p}-\mathbf v_\star$, position and velocity of the planet relative to the star;
  \item $\mathbf s$, spin vector of the planet;
  \item $M^\star$, $M_{\mathrm p}$, $R$, and the gravitational constant $\G$;
  \item the Love-number spectra $\texttt{rek2}(\chi)$, $\texttt{imk2}(\chi)$ of the planet (Section~\ref{sec:love});
  \item the history of the accepted steps, from which the mean rates $\bar n$, $\dot{\bar\lambda}$, $\dot{\bar\Omega}$ are computed (Section~\ref{sec:rates}), and the mean motions of the other orbiting bodies, which fix the averaging window.
\end{itemize}

\paragraph{Outputs.}
\begin{itemize}
  \item The force of the star on the planet, $\Fsp=M^\star\nabla T|_{\vr^\star}$, in inertial axes (convention~\eqref{eq:force-convention}).
  \item The torque on the spin, $\boldsymbol\Gamma=\dd(I\mathbf s)/\dd t$ (Section~\ref{sec:torque}).
\end{itemize}

\paragraph{How the force is applied.} The tidal interaction is an internal force of the star--planet pair:
\begin{itemize}
  \item the planet receives $+\Fsp/M_{\mathrm p}$ and the star $-\Fsp/M^\star$, in the barycentric frame;
  \item the relative acceleration is then $\Fsp(1/M_{\mathrm p}+1/M^\star)$ automatically;
  \item in a multi-planet system, the reaction on the star produces the indirect terms in the motion of the other planets.
\end{itemize}
No extra factor and no indirect term must be added by hand.

\section{Geometry from the state vectors}\label{sec:geometry}
All geometric quantities are computed from $(\vr,\mathbf v,\mathbf s)$. No Euler angle of the spin axis and no inertial reference direction is needed.

\paragraph{Orbit of the star seen from the planet.} The star is at $\vr^\star=-\vr$, with velocity $\mathbf v^\star=-\mathbf v$. Its orbit has the angular momentum per unit reduced mass $\mathbf h=\vr^\star\times\mathbf v^\star=\vr\times\mathbf v$, the same as that of the planet. With $\mu_{\rm orb}=\G(M^\star+M_{\mathrm p})$ and $r^\star=\|\vr^\star\|$,
\begin{equation}
  a=\Big(\frac{2}{r^\star}-\frac{\|\mathbf v\|^2}{\mu_{\rm orb}}\Big)^{-1},\qquad
  \mathbf e_{\rm v}=\frac{\mathbf v^\star\times\mathbf h}{\mu_{\rm orb}}-\frac{\vr^\star}{r^\star},\qquad
  e=\|\mathbf e_{\rm v}\|,\qquad \bar n_{\rm osc}=\sqrt{\mu_{\rm orb}/a^3}. \label{eq:elements-a-e}
\end{equation}

\paragraph{Nodal frame.} Let $\hat{\mathbf s}=\mathbf s/\|\mathbf s\|$ be the spin axis, $\Op=\|\mathbf s\|$ the spin rate, and $\hat{\mathbf h}=\mathbf h/\|\mathbf h\|$ the orbit normal. The node of the orbit on the equator is the intersection of the two planes, i.e.\ the direction $\hat{\mathbf s}\times\hat{\mathbf h}$, whose norm is $\sin i$. The frame is built in three steps.
\begin{enumerate}
  \item \textbf{Direction of the node.}
  \[
    \mathbf N_0=\begin{cases}\hat{\mathbf s}\times\hat{\mathbf h}, & \text{if}\ \|\hat{\mathbf s}\times\hat{\mathbf h}\|\ge\varepsilon_N,\\[1mm] \hat{\mathbf s}\times\mathbf e_k, & \text{otherwise}.\end{cases}
  \]
  $\mathbf e_k$ is the inertial axis with the smallest $|\hat s_k|$, and $\varepsilon_N=10^{-10}$.
  \item \textbf{Projection onto the equator and normalisation.}
  \[
    \mathbf N_\perp=\mathbf N_0-(\mathbf N_0\cdot\hat{\mathbf s})\,\hat{\mathbf s},\qquad
    \hat{\mathbf N}=\frac{\mathbf N_\perp}{\|\mathbf N_\perp\|}.
  \]
  \item \textbf{Third axis.}
  \[
    \hat{\mathbf Y}=\hat{\mathbf s}\times\hat{\mathbf N}.
  \]
\end{enumerate}
In compact form,
\begin{equation}
  \hat{\mathbf N}=\frac{\mathbf N_0-(\mathbf N_0\cdot\hat{\mathbf s})\,\hat{\mathbf s}}{\big\|\mathbf N_0-(\mathbf N_0\cdot\hat{\mathbf s})\,\hat{\mathbf s}\big\|},\qquad
  \hat{\mathbf Y}=\hat{\mathbf s}\times\hat{\mathbf N}. \label{eq:node-frame}
\end{equation}
\begin{itemize}
  \item $(\hat{\mathbf N},\hat{\mathbf Y},\hat{\mathbf s})$ is the nodal frame of Chapter~\ref{ch:second}: $\hat{\mathbf N}$ points to the ascending node of the star's orbit on the equator.
  \item \textbf{The denominator.} In the general case, $\mathbf N_0=\hat{\mathbf s}\times\hat{\mathbf h}$ is already orthogonal to $\hat{\mathbf s}$, so $\mathbf N_\perp=\mathbf N_0$ and $\|\mathbf N_\perp\|=\|\hat{\mathbf s}\times\hat{\mathbf h}\|=\sin i\ge\varepsilon_N$. In the other case, $\|\mathbf N_\perp\|=\|\hat{\mathbf s}\times\mathbf e_k\|=\sqrt{1-\hat s_k^2}\ge\sqrt{2/3}$, because the smallest component of a unit vector satisfies $\hat s_k^2\le1/3$. The denominator is therefore never zero.
  \item \textbf{The projection.} Mathematically $\mathbf N_\perp=\mathbf N_0$. The projection only removes the rounding error, so that $\hat{\mathbf N}\cdot\hat{\mathbf s}=0$ to machine precision.
  \item \textbf{The threshold.} When $\hat{\mathbf s}\parallel\hat{\mathbf h}$ ($i=0$ or $\pi$), the node is undefined and any equatorial direction can be used: the force does not depend on it (Section~\ref{sec:validation}, test V7). The threshold $\varepsilon_N$ only avoids dividing by a vanishing $\sin i$; it does not switch to another formula for the force.
\end{itemize}

\paragraph{Inclination and angles in the orbital plane.}
\begin{equation}
  i=\operatorname{atan2}\big(\|\hat{\mathbf s}\times\hat{\mathbf h}\|,\ \hat{\mathbf h}\cdot\hat{\mathbf s}\big)\in[0,\pi],\qquad
  \hat{\mathbf W}=\frac{\hat{\mathbf h}\times\hat{\mathbf N}}{\|\hat{\mathbf h}\times\hat{\mathbf N}\|},\qquad
  \hat{\mathbf N}'=\hat{\mathbf W}\times\hat{\mathbf h}. \label{eq:inclination}
\end{equation}
$\hat{\mathbf N}'$ and $\hat{\mathbf W}$ span the orbital plane, with $\hat{\mathbf N}'=\hat{\mathbf N}$ when the node is defined. A retrograde rotation corresponds to $i>\pi/2$; no sign convention is needed for $\Op$.

\paragraph{Elements referred to the equator.} With $\hat{\vr}^\star=\vr^\star/r^\star$,
\begin{gather}
  u=\operatorname{atan2}\big(\hat{\vr}^\star\cdot\hat{\mathbf W},\ \hat{\vr}^\star\cdot\hat{\mathbf N}'\big),\qquad
  \omega=\operatorname{atan2}\big(\mathbf e_{\rm v}\cdot\hat{\mathbf W},\ \mathbf e_{\rm v}\cdot\hat{\mathbf N}'\big),\qquad f=u-\omega, \label{eq:elements-angles}\\
  E=\operatorname{atan2}\big(\sqrt{1-e^2}\,\sin f,\ e+\cos f\big),\qquad \M=E-e\sin E,\qquad u_\M\equiv\omega+\M. \notag
\end{gather}
\begin{itemize}
  \item $\omega$ is the argument of pericentre of the \emph{star's} orbit, referred to the planet's equator. It equals that of the planet plus $\pi$.
  \item $u_\M$ is the argument of the mean latitude. It is defined even when $e=0$: the force depends on $\omega$ and $\M$ only through $u_\M$ and through $\M$ multiplied by $G_{2pq}(e)=\mathcal O(e^{|q|})$. If $e=0$, set $\omega=0$.
\end{itemize}

\paragraph{Position of the star in the nodal frame.}
\begin{equation}
  \rho\equiv\hat{\vr}^\star\cdot\hat{\mathbf N}+\I\,\hat{\vr}^\star\cdot\hat{\mathbf Y}=\sin\theta^\star e^{\I\lambda_N},\qquad
  c^\star\equiv\cos\theta^\star=\hat{\vr}^\star\cdot\hat{\mathbf s}, \label{eq:rho}
\end{equation}
where $\theta^\star$ is the colatitude of the star and $\lambda_N$ its longitude counted from the node, as in~\eqref{eq:star-geometry}.

\section{From the gradient to the implementable form}\label{sec:implementable}
We start from the Cartesian gradient~\eqref{eq:grad-cart} at $l=2$, evaluated at the star. Three steps make it implementable: the change to the nodal frame, the elimination of $\Omega-\thetap$, and the polynomial form of the angular functions.

\paragraph{Notation for $l=2$.} With the elements of Section~\ref{sec:geometry},
\begin{equation}
  \mathcal K_{2mpq}=\G M^\star\frac{R^5}{a^3}\,c_m\,F_{2mp}(i)\,G_{2pq}(e),\qquad c_0=1,\ c_1=\tfrac13,\ c_2=\tfrac1{12}, \label{eq:K2}
\end{equation}
from~\eqref{eq:Klmpq}, and $\kb_2\equiv\kb_2(\bar\sigma_{2mpq})$ (Section~\ref{sec:love}). Only the terms with $F_{2mp}(i)\neq0$ and $|q|\le q_{\max}(e)$ (Section~\ref{sec:truncation}) are kept.

\paragraph{Step 1: phases at the star.} By~\eqref{eq:star-geometry}, the body-fixed longitude of the star is $\varphi^\star=(\Omega-\thetap)+\lambda_N$. In the phase of~\eqref{eq:grad-cart}, $\Lambda_{2mpq}=\zeta_{2mpq}-m\varphi^\star$, the term $m(\Omega-\thetap)$ of $\nu_{2mpq}$ cancels:
\begin{equation}
  \Lambda_{2mpq}=\Psi_{mpq}-m\lambda_N,\qquad \Psi_{mpq}\equiv(2-2p)\,\omega+(2-2p+q)\,\M+f_{2m}=(2-2p)\,u_\M+q\,\M+f_{2m}. \label{eq:Psi}
\end{equation}
The second form uses $(2-2p)\omega+(2-2p+q)\M=(2-2p)(\omega+\M)+q\M$ and is regular as $e\to0$.

\paragraph{Step 2: phase factor.} With $|\kb_2|e^{-\I\epsilon_2}=\kb_2$, the factor $|\kb_2|e^{\pm\I X}$ of~\eqref{eq:grad-cart}, $X=\Lambda-\epsilon_2$, becomes
\begin{equation}
  |\kb_2|\,e^{\I X}=\mathcal E_{mpq}\,e^{-\I m\lambda_N},\qquad \mathcal E_{mpq}\equiv\kb_2(\bar\sigma_{2mpq})\,e^{\I\Psi_{mpq}}, \label{eq:Ecal}
\end{equation}
and $|\kb_2|e^{-\I X}=\mathcal E^*_{mpq}e^{\I m\lambda_N}$.

\paragraph{Step 3: angular functions.} The degree-3 functions factor as
\[
  \PK_{3m'}(\cos\theta)=\sin^{m'}\theta\;\mathcal P_{3m'}(\cos\theta),\qquad \mathcal P_{3m'}=\frac{\dd^{m'}P_3}{\dd x^{m'}}
\]
(Appendix~\ref{app:harmonics}), with
\begin{equation}
  \mathcal P_{30}(x)=\tfrac12(5x^3-3x),\qquad \mathcal P_{31}(x)=\tfrac32(5x^2-1),\qquad \mathcal P_{32}(x)=15x,\qquad \mathcal P_{33}(x)=15. \label{eq:P3m}
\end{equation}
Hence, at the star, by~\eqref{eq:rho},
\[
  \PK_{3m'}(c^\star)\,e^{\pm\I m'\lambda_N}=\rho^{m'}\mathcal P_{3m'}(c^\star)\quad(+),\qquad \bar\rho^{\,m'}\mathcal P_{3m'}(c^\star)\quad(-).
\]
\begin{itemize}
  \item $Z$-component: $|\kb_2|\,\PK_{3m}\cos X=\Re\big[\PK_{3m}\,e^{-\I m\lambda_N}\mathcal E_{mpq}\big]=\mathcal P_{3m}\,\Re\big[\bar\rho^{\,m}\mathcal E_{mpq}\big]$.
  \item Order $m+1$: $|\kb_2|\,e^{\I\lambda_N}\PK_{3,m+1}e^{-\I X}=\PK_{3,m+1}\,e^{\I(m+1)\lambda_N}\mathcal E^*_{mpq}=\rho^{m+1}\mathcal P_{3,m+1}\,\mathcal E^*_{mpq}$.
  \item Order $m-1$: $|\kb_2|\,e^{\I\lambda_N}\PK_{3,m-1}e^{\I X}=\PK_{3,m-1}\,e^{-\I(m-1)\lambda_N}\mathcal E_{mpq}=\bar\rho^{\,m-1}\mathcal P_{3,m-1}\,\mathcal E_{mpq}$.
\end{itemize}
In the last two lines, $e^{\I\varphi^\star}=e^{\I(\Omega-\thetap)}e^{\I\lambda_N}$ was used; the factor $e^{\I(\Omega-\thetap)}$ is removed by the rotation~\eqref{eq:nodal-rot} to the nodal frame. Steps 1--3 were checked numerically term by term.

\paragraph{Result.} The gradient of $T$ at the star, in the nodal frame, is
\begin{equation}
  \boxed{\begin{aligned}
  V_Z&=\frac{1}{(r^\star)^4}\sum_{m,p,q}(3-m)\,\mathcal K_{2mpq}\,\mathcal P_{3m}(c^\star)\,\Re\big[\bar\rho^{\,m}\mathcal E_{mpq}\big],\\
  V_X+\I V_Y&=\frac{1}{(r^\star)^4}\sum_{p,q}\mathcal K_{20pq}\,\mathcal P_{31}(c^\star)\,\rho\,\Re\,\mathcal E_{0pq}\\
  &\quad+\frac{1}{2(r^\star)^4}\sum_{m=1}^{2}\sum_{p,q}\mathcal K_{2mpq}\Big[\mathcal P_{3,m+1}(c^\star)\,\rho^{m+1}\mathcal E^*_{mpq}\\
  &\hspace{9em}-(3-m)(4-m)\,\mathcal P_{3,m-1}(c^\star)\,\bar\rho^{\,m-1}\mathcal E_{mpq}\Big].
  \end{aligned}} \label{eq:impl-single}
\end{equation}
where:
\begin{itemize}
  \item $V_Z$ is the component of $\nabla T$ along the spin axis $\hat{\mathbf s}$, and $V_X+\I V_Y$ its components along $\hat{\mathbf N}$ and $\hat{\mathbf Y}$ (nodal frame~\eqref{eq:node-frame});
  \item $r^\star$ is the distance star--planet, $\rho=\sin\theta^\star e^{\I\lambda_N}$ and $c^\star=\cos\theta^\star$ give the direction of the star in the nodal frame~\eqref{eq:rho};
  \item $\mathcal K_{2mpq}$ is the real coefficient~\eqref{eq:K2}: masses, radius, $a$, $F_{2mp}(i)$, $G_{2pq}(e)$;
  \item $\mathcal P_{3m'}$ are the polynomials~\eqref{eq:P3m};
  \item $\mathcal E_{mpq}=\kb_2(\bar\sigma_{2mpq})\,e^{\I\Psi_{mpq}}$ combines the Love number of the term and its phase $\Psi_{mpq}=(2-2p)u_\M+q\M+f_{2m}$ \eqref{eq:Psi}--\eqref{eq:Ecal};
  \item the sums run over $m,p\in\{0,1,2\}$ and $|q|\le q_{\max}(e)$ (Section~\ref{sec:truncation}).
\end{itemize}
Properties:
\begin{itemize}
  \item It holds for any obliquity $0\le i\le\pi$. No separate formula is needed at $i=0$. A faster form for a coplanar orbit is given in Section~\ref{sec:coplanar-nbody}, as an optimisation only.
  \item It contains no division by $\sin i$, $\sin\theta^\star$ or $e$. The angular functions are polynomials in $(\rho,\bar\rho,c^\star)$.
  \item The exact geometry is used: $r^\star$ and the direction of the star come from the state vectors, not from the series. The truncation acts only on the Kaula expansion of the tidal-raising potential.
\end{itemize}
The double transformation~\eqref{eq:grad-node}, which uses the elements only, is given in Appendix~\ref{app:double} for reference.

\section{Tidal frequencies and Love numbers}\label{sec:love}
\subsection{Frequencies with mean rates}\label{sec:rates}
The Love number is defined for a harmonic forcing of fixed frequency (H1). Each Kaula term is quasi-harmonic, with a slowly varying rate (H3). The frequency used in $\kb_2$ is therefore the rate of $\nu_{2mpq}$ averaged over the \emph{fast} angles: the orbital period and the synodic periods with the other planets. The slow angles (secular precessions, resonant angles, super-periods) are not averaged; they are part of the forcing frequency. The phases $\Psi_{mpq}$ always use the osculating $u_\M$ and $\M$ of the current instant.

\paragraph{Regular form.} The rate~\eqref{eq:sigma} of Kaula's angle can be written in three equivalent forms, derived in Appendix~\ref{app:frequency}:
\begin{itemize}
  \item Kaula's form: $\sigma_{2mpq}=(2-2p+q)\,n-m\,\Op+(2-2p)\,\dot\omega+m\,\dot\Omega$;
  \item with the longitude of pericentre $\varpi=\Omega+\omega$: $\sigma_{2mpq}=(2-2p+q)\,n-m\,\Op+(2-2p)\,\dot\varpi+(m-2+2p)\,\dot\Omega$. In the coplanar case only $k=m-2+2p=0$ occurs and this is the usual form $(2-2p+q)n-m\Op+(2-2p)(\dot\Omega+\dot\omega)$;
  \item with the mean longitude $\lambda=\Omega+\omega+\M$, the form used here.
\end{itemize}
With $\lambda\equiv\Omega+\omega+\M$, the mean longitude of the star counted in the equator, and $k\equiv m-2+2p$, the rate at $l=2$ reads, identically,
\begin{equation}
  \boxed{\bar\sigma_{2mpq}=(2-2p)\,\dot{\bar\lambda}+q\,\bar n+k\,\dot{\bar\Omega}-m\,\Op,\qquad k=m-2+2p.} \label{eq:sigma-regular}
\end{equation}
where (a bar denotes the average over the fast angles, defined below):
\begin{itemize}
  \item $\dot{\bar\lambda}$ is the mean rate of the mean longitude $\lambda$ of the star; it is close to $n$ and contains the precession of the pericentre and of the node;
  \item $\bar n$ is the mean motion $\sqrt{\G(M^\star+M_{\mathrm p})/a^3}$ averaged over the window; it replaces $\dot\M$;
  \item $\dot{\bar\Omega}$ is the mean precession rate of the node of the star's orbit on the planet's equator, zero for an isolated Kepler orbit;
  \item $\Op=\|\mathbf s\|$ is the spin rate of the planet;
  \item $m,p,q$ are the indices of the term.
\end{itemize}
Properties:
\begin{itemize}
  \item $\dot{\bar\lambda}$ stays defined as $e\to0$ and as $i\to0$. The osculating $\dot\omega$ alone behaves like $1/e$, and $\dot\Omega$ like $1/\sin i$.
  \item $\dot{\bar\Omega}$ appears only with $k\neq0$. These terms have $F_{2mp}(i)=\mathcal O(i^{|k|})$ and vanish as $i\to0$, where $\Omega$ is undefined. For $m=2,p=0$ and $m=0,p=1$ ($k=0$, the only terms of the coplanar case), the frequency depends on $\dot{\bar\lambda}$, $\bar n$ and $\Op$ only.
\end{itemize}

\paragraph{Averaging window.} The window $\Delta T$ must remove the fast oscillations and keep the slow ones:
\begin{itemize}
  \item longer than the orbital period $P_i=2\pi/n_i$ of the planet, and than its synodic periods with the other bodies, $P_{{\rm syn},ij}=2\pi/|n_i-n_j|$, which are the periods of the conjunctions;
  \item much shorter than the slow time scales (secular precessions, libration of a resonance), which are part of the forcing.
\end{itemize}
It is computed by the code, at each accepted step, from the osculating mean motions $n_j$ of all the other orbiting bodies:
\begin{equation}
  \boxed{\Delta T_i=\min\Big(\max\big(\max_{j\neq i}P_{{\rm syn},ij},\ P_i\big),\ 10\,P_i\Big),\qquad P_{{\rm syn},ij}=\frac{2\pi}{|n_i-n_j|},\quad P_i=\frac{2\pi}{n_i}.} \label{eq:window}
\end{equation}
\begin{itemize}
  \item The longest synodic period covers all the conjunction frequencies of the planet.
  \item The lower bound $P_i$ applies to a single planet, the upper bound $10P_i$ to two planets with very close periods, whose synodic period would otherwise become longer than the slow time scales.
  \item The user gives no parameter.
\end{itemize}

\paragraph{Estimation in a run.} Keep, for each tidally active planet, a history of the accepted steps $(t_j,\lambda_j,\Omega_j,n_{{\rm osc},j})$ covering at least $\Delta T$:
\begin{enumerate}
  \item $\Omega_j$ is accumulated from the rotation of $\hat{\mathbf N}$ about $\hat{\mathbf s}$: $\Delta\Omega=\operatorname{atan2}\big(\hat{\mathbf s}\cdot(\hat{\mathbf N}_{j-1}\times\hat{\mathbf N}_j),\ \hat{\mathbf N}_{j-1}\cdot\hat{\mathbf N}_j\big)$;
  \item $\lambda_j=\lambda_{j-1}+\operatorname{wrap}\big(\Delta\Omega+u_{\M,j}-u_{\M,j-1}\big)$, with $\operatorname{wrap}$ to $(-\pi,\pi]$;
  \item $\dot{\bar\lambda}=[\lambda(t)-\lambda(t-\Delta T)]/\Delta T$ and $\dot{\bar\Omega}=[\Omega(t)-\Omega(t-\Delta T)]/\Delta T$, the values at $t-\Delta T$ being interpolated linearly between the two stored samples around it; $\bar n$ is the mean of $n_{\rm osc}$ over the window;
  \item before one window has elapsed: $\dot{\bar\lambda}=\bar n=n_{\rm osc}$, $\dot{\bar\Omega}=0$.
\end{enumerate}

\paragraph{Accuracy.}
\begin{itemize}
  \item Two planets near the 3:2 resonance: the rule~\eqref{eq:window} gives one synodic period. One integration step ($\Delta T=P/20$) gives the osculating rate, with an error of $4.5\times10^{-4}n$ on $\bar\sigma_{2200}$ and up to $45\,\%$ on the force near synchronism. One synodic period gives $9\times10^{-5}n$ and $4\,\%$.
  \item TRAPPIST-1 b, c, d, e, with their resonant chain: the rule gives $2.0$ to $3.0\,P_i$. The error on $\bar\sigma_{2200}$ is $2\times10^{-5}n$ to $6\times10^{-5}n$, against $5\times10^{-4}n$ to $8\times10^{-4}n$ for one step and $2\times10^{-4}n$ for one orbital period. Over the four planets, the worst error of the rule is within $10\,\%$ of the best of the fixed windows tested ($P_i$ to $10P_i$). No single window is the best for every planet.
  \item The history is updated on \emph{accepted} steps only. Intermediate force evaluations (sub-steps, predictor--corrector iterations) only read it, so that the rates do not depend on the trial states of the integrator.
\end{itemize}

\subsection{Love numbers from tabulated spectra}
The spectra $\texttt{rek2}(\chi)$ and $\texttt{imk2}(\chi)\ge0$ are given at physical frequencies $\chi=|\sigma|\ge0$, with $\texttt{imk2}>0$ for a lag. By~\eqref{eq:love-spectra},
\begin{equation}
  \kb_2(\bar\sigma)=\texttt{rek2}(|\bar\sigma|)-\I\,\sgn(\bar\sigma)\,\texttt{imk2}(|\bar\sigma|),\qquad \kb_2(0)=\texttt{rek2}(0). \label{eq:love-impl}
\end{equation}

\paragraph{Preparation of the table.} Keep the strictly increasing prefix $\chi_1<\dots<\chi_N$ of the tabulated frequencies: it stops at a padding that repeats the last row or that falls back to zero. Inside $[\chi_1,\chi_N]$, interpolate linearly.

\paragraph{Extrapolation above $\chi_N$.} Real rheologies have power-law tails at high frequency: $\texttt{imk2}\propto\chi^{-\alpha}$ and $\texttt{rek2}\to k_\infty+C\chi^{b}$, with $b=-\alpha$ (Andrade) or $b=-2\alpha$ (Maxwell, $\alpha=1$). With the last two (three) distinct points:
\begin{equation}
  \begin{gathered}
  \beta=\frac{\ln(\texttt{imk2}_N/\texttt{imk2}_{N-1})}{\ln(\chi_N/\chi_{N-1})}\ \text{clamped to}\ [-1,0],\qquad
  \texttt{imk2}(\chi)=\texttt{imk2}_N\Big(\frac{\chi}{\chi_N}\Big)^{\beta},\\
  \texttt{rek2}(\chi)=\texttt{rek2}_N+(\texttt{rek2}_N-\texttt{rek2}_{N-1})\,\frac{\chi^{b}-\chi_N^{b}}{\chi_N^{b}-\chi_{N-1}^{b}},
  \end{gathered} \label{eq:extrap}
\end{equation}
where $\chi_N>\chi_{N-1}>\chi_{N-2}$ are the last three distinct frequencies of the table, $\texttt{rek2}_j$ and $\texttt{imk2}_j$ the tabulated values there, $\beta$ the high-frequency slope of $\texttt{imk2}$ in log--log, and $b$ solves $\dfrac{\chi_N^b-\chi_{N-1}^b}{\chi_{N-1}^b-\chi_{N-2}^b}=\dfrac{\texttt{rek2}_N-\texttt{rek2}_{N-1}}{\texttt{rek2}_{N-1}-\texttt{rek2}_{N-2}}$ (a fit of $\texttt{rek2}=k_\infty+C\chi^b$ through the last three points), clamped between $2\beta$ and $\beta$, i.e.\ between the Maxwell and the Andrade behaviours.
\begin{itemize}
  \item If $\texttt{imk2}_N$ or $\texttt{imk2}_{N-1}$ is not positive, or $\beta\ge0$, both values are held constant. If $\texttt{rek2}$ is not monotonic over the last three points, $b=\beta$ is used.
  \item For $\beta<0$, $\chi^b\to0$: $\texttt{rek2}$ tends to a finite limit and $\texttt{imk2}$ to zero. The extrapolated values stay finite and $\texttt{imk2}\ge0$ for any $\chi$.
\end{itemize}

\paragraph{Extrapolation below $\chi_1$} (only if $\chi_1>0$). The same rule is applied with the first points, $\beta$ clamped to $[0,1]$ ($\texttt{imk2}\propto\chi^{\beta}$ below the dissipation peak), and $\texttt{imk2}(0)=0$.

\paragraph{Accuracy.} Tested on spectra truncated by one decade and compared with the exact values:
\begin{center}\small
\begin{tabular}{lcc}\toprule
Spectrum (cut one decade before its end) & error on \texttt{rek2} & error on \texttt{imk2}\\\midrule
Maxwell, homogeneous body & $3\times10^{-14}$ & $4\times10^{-9}$\\
Andrade, $\alpha=0.2$ / $0.3$ & $3\times10^{-5}$ / $1\times10^{-5}$ & $1.3\times10^{-3}$\\
Burgers & $6\times10^{-5}$ & $4\times10^{-2}$\\
TRAPPIST-1 b--h spectra (Revol et al.\ 2024) & $\le5\times10^{-3}$ & $\le8\times10^{-3}$\\
Low-frequency end, Maxwell / Andrade (cut below the peak) & $\le2\times10^{-7}$ & $\le5\times10^{-3}$\\
Holding the end values (no extrapolation), same cases & up to $7\times10^{-2}$ & $0.6$ to $9$\\\bottomrule
\end{tabular}
\end{center}
No rule can predict a feature that lies outside the data, such as a second relaxation peak beyond $\chi_N$: the spectra should cover the forcing frequencies of the run, with the extrapolation as a safeguard.

\section{Truncation and eccentricity functions}\label{sec:truncation}
The force must be accurate for any Love-number spectrum, since the spectrum can vary strongly with frequency. In the worst case, each term may have an arbitrary $\kb_{2,mpq}$ with $|\kb_{2,mpq}|\le\kappa_{\max}$. The force of a term is linear in $(\Re\kb,\Im\kb)$, $\vF_{mpq}=\mathsf L_{mpq}(\Re\kb,\Im\kb)^{\mathsf T}$, so the truncation error is bounded by
\[
  \|\delta\vF\|\le\kappa_{\max}\sum_{|q|>q_{\max}}\sum_{m,p}\|\mathsf L_{mpq}\|_2.
\]
The table below keeps $|q|\le q_{\max}(e)$. It guarantees, for any spectrum with $|\kb_2|\le\kappa_{\max}$,
\begin{equation}
  \|\delta\vF\|\le10^{-5}\,\kappa_{\max}F_0(r^\star),\qquad \|\delta\boldsymbol\Gamma\|\le10^{-5}\,\kappa_{\max}\,r^\star F_0(r^\star),\qquad F_0(r^\star)=\frac{3\G M^{\star2}R^5}{(r^\star)^7}, \label{eq:guarantee}
\end{equation}
at every point of the orbit, where $F_0(r^\star)$ is the instantaneous elastic force per unit Love number:
\begin{center}
\renewcommand{\arraystretch}{1.2}\setlength{\tabcolsep}{4pt}\small
\begin{tabular}{l|cccccccccccc}\toprule
$e\le$ & 0.001 & 0.01 & 0.02 & 0.05 & 0.10 & 0.15 & 0.20 & 0.25 & 0.30 & 0.35 & 0.40 & 0.50\\\midrule
$q_{\max}$ & 2 & 3 & 3 & 5 & 7 & 10 & 12 & 15 & 18 & 23 & 28 & 42\\\bottomrule
\end{tabular}
\end{center}
\begin{itemize}
  \item Each value is the worst over the points of a 35-point grid in $e$ that fall in the bin. For each $e$ there are 714 states: 7 inclinations from 0 to $\pi$, 3 arguments of pericentre, and 34 positions on the orbit (32 uniform in $\M$, plus apocentre and pericentre). The worst state is always near the apocentre, where the force is weakest. The bound was checked again on independent states.
  \item The torque bound, divided by $r^\star F_0$, never exceeds the force bound divided by $F_0$, since $\|\vr^\star\times\delta\vF\|\le r^\star\|\delta\vF\|$.
  \item The guarantee is relative to $\kappa_{\max}F_0$, a scale set mostly by the elastic part, while the secular evolution is driven by the dissipative part $\Im\kb_2$. The tolerance $10^{-5}$ resolves the dissipative part to $10^{-3}$ in the worst case when $|\Im\kb_2|/\kappa_{\max}\gtrsim10^{-2}$; in general, choose a tolerance of order $10^{-3}|\Im\kb_2|/\kappa_{\max}$. The tables for tolerances $10^{-2}$ to $10^{-6}$ are in Appendix~\ref{app:truncation}.
  \item The table covers $e\le0.5$. Beyond, $q_{\max}$ must be extended with the same method.
\end{itemize}

\paragraph{Eccentricity functions.} $G_{2pq}(e)$ is needed up to $|q|=42$. The power series truncated at $e^{7}$ is not sufficient: its error reaches $8\times10^{-4}$ at $e=0.2$ and $2\times10^{-2}$ at $e=0.3$. $G_{2pq}$ is computed exactly with Tisserand's formula (Appendix~\ref{app:kaula}, eq.~\eqref{eq:tisserand}), with at most 22 terms for $e\le0.5$ and an absolute error below $10^{-13}$. The symmetries $G_{22q}=G_{20,-q}$ and $G_{21q}=G_{21,-q}$ halve the work. The inclination functions $F_{2mp}(i)$ are the closed forms of Table~\ref{tab:F2mp}.

\paragraph{Fast evaluation of $G_{2pq}$.} Tisserand's formula costs up to $10^{-4}$\,s per force evaluation at $e=0.4$, and a cache keyed on $e$ does not help: the eccentricity rebuilt from the state vectors changes in its last bits at every call. For $e\le0.5$ the functions are therefore read from Chebyshev tables of their smooth part,
\begin{equation}
  G_{2pq}(e)=\beta_e^{|q|}\,H_{pq}(u),\qquad u=\beta_e^2,\qquad \beta_e=\frac{e}{1+\sqrt{1-e^2}},\qquad
  H_{pq}(u)\simeq\sum_{k=0}^{8}c^{(j)}_{pqk}\,T_k\big(t_j(u)\big), \label{eq:G-table}
\end{equation}
where:
\begin{itemize}
  \item $u\in[0,u_{\max}]$, $u_{\max}=\beta_e^2(0.5)=0.0718$, is cut into 8 equal intervals $j$; $t_j(u)\in[-1,1]$ is the local variable of interval $j$ and $T_k$ are the Chebyshev polynomials (Clenshaw's recurrence);
  \item $H_{pq}=G_{2pq}/\beta_e^{|q|}$ is a smooth function of $u$, since $G_{2pq}=\mathcal O(e^{|q|})$; the factor $\beta_e^{|q|}$ is applied afterwards;
  \item the coefficients $c^{(j)}_{pqk}$ are computed from Tisserand's formula (45 terms) at the 9 Chebyshev nodes of each interval, once per run (about 0.1\,s).
\end{itemize}
The absolute error is $3\times10^{-15}$ for $e\le0.5$ and $|q|\le q_{\max}(e)$. Above $e=0.5$, Tisserand's formula is used directly.

\section{Force and torque in inertial axes}\label{sec:torque}
\paragraph{Force.} From~\eqref{eq:impl-single} and the convention~\eqref{eq:force-convention},
\begin{equation}
  \boxed{\Fsp=M^\star\Big[\Re(V_X+\I V_Y)\,\hat{\mathbf N}+\Im(V_X+\I V_Y)\,\hat{\mathbf Y}+V_Z\,\hat{\mathbf s}\Big].} \label{eq:impl-force}
\end{equation}
where $\Fsp$ is the force of the star on the planet in inertial axes, $M^\star$ the mass of the star, $V_X,V_Y,V_Z$ the components~\eqref{eq:impl-single} of $\nabla T$ at the star, and $(\hat{\mathbf N},\hat{\mathbf Y},\hat{\mathbf s})$ the nodal frame expressed in inertial axes~\eqref{eq:node-frame}.

\paragraph{Torque.} Neglect the tide in the star. The only internal interaction of the pair is $\Fsp=-\Fps$, applied at the centres of mass. With the reduced mass $\mu=M^\star M_{\mathrm p}/(M^\star+M_{\mathrm p})$, the orbital angular momentum $\mathbf L_{\rm orb}=\mu\,\vr\times\dot\vr$ evolves as $\dot{\mathbf L}_{\rm orb}=\vr\times\Fsp$, because gravity is central. The conservation of $\mathbf L_{\rm orb}+\mathbf L_{\rm spin}$ then imposes
\begin{equation}
  \boxed{\frac{\dd(I\mathbf s)}{\dd t}=\boldsymbol\Gamma=-\,\vr\times\Fsp=M^\star\,\vr^\star\times\nabla T\big|_{\vr^\star}.} \label{eq:torque}
\end{equation}
where $I$ is the moment of inertia of the planet, $\mathbf s$ its spin vector, $\vr=\vr_{\mathrm p}-\vr_\star$ the position of the planet relative to the star, $\vr^\star=-\vr$ the position of the star seen from the planet, and $\Fsp$ the force~\eqref{eq:impl-force}.
\begin{itemize}
  \item The torque must be computed with \emph{the same instantaneous} force as the one applied to the orbit. A torque built from an orbit-averaged force exchanges, at each step, a different angular momentum from the one taken from the orbit, and $\mathbf L_{\rm tot}$ drifts.
  \item The component along the spin axis, $\boldsymbol\Gamma\cdot\hat{\mathbf s}$, changes the spin rate; the perpendicular components change the direction of the spin axis.
  \item The constant-time-lag model \citep{Hut1981,Bolmont2015} satisfies~\eqref{eq:torque} exactly.
\end{itemize}

\section{Coplanar orbit: a faster form}\label{sec:coplanar-nbody}
When the orbit lies in the equator of the planet ($i=0$), only $(m,p)=(0,1)$ and $(2,0)$ survive in~\eqref{eq:impl-single}, since $F_{201}(0)=-\tfrac12$, $F_{220}(0)=3$ and the other $F_{2mp}(0)$ vanish. The star is in the equator: $c^\star=0$, $|\rho|=1$, $\mathcal P_{31}(0)=-\tfrac32$, $\mathcal P_{33}(0)=15$ and $\mathcal P_{30}(0)=\mathcal P_{32}(0)=0$. Hence $V_Z=0$, and $V_X+\I V_Y=\rho\,(w_r+\I w_\varphi)$ with real $w_r$, $w_\varphi$. The gradient has a radial and an azimuthal part only:
\begin{equation}
\boxed{\begin{gathered}
  \nabla T\big|_{\vr^\star}=\frac{\G M^\star R^5}{a^3(r^\star)^4}\big(w_r\,\hat\vr^\star+w_\varphi\,\hat{\boldsymbol\varphi}\big),\qquad \hat{\boldsymbol\varphi}=\hat{\mathbf s}\times\hat\vr^\star,\\
  w_r=\sum_{|q|\le q_{\max}}\Big[\tfrac34\,G_{21q}(e)\,\Re\big(\kb_{01q}\,e^{\I q\M}\big)+\tfrac94\,G_{20q}(e)\,\Re Z_q\Big],\\
  w_\varphi=-\tfrac32\sum_{|q|\le q_{\max}}G_{20q}(e)\,\Im Z_q,\qquad Z_q=\kb_{20q}\,e^{\I[(2+q)\M-2f]},
\end{gathered}} \label{eq:coplanar-nbody}
\end{equation}
where:
\begin{itemize}
  \item $\kb_{01q}=\kb_2(q\bar n)$ and $\kb_{20q}=\kb_2(2\dot{\bar\lambda}+q\bar n-2\Op)$ are the Love numbers at the frequencies~\eqref{eq:sigma-regular} with $i=0$;
  \item $f$ is the true anomaly of the star: $\cos f=\hat{\mathbf e}_{\rm v}\cdot\hat\vr^\star$, $\sin f=\hat{\mathbf h}\cdot(\hat{\mathbf e}_{\rm v}\times\hat\vr^\star)$, with $\hat{\mathbf e}_{\rm v}=\mathbf e_{\rm v}/e$; if $e=0$, $f=\M=0$;
  \item only $a$, $e$, $r^\star$, $\hat\vr^\star$, $\M$ and $f$ are needed: no node, no $\omega$, no $F_{2mp}$ and no $\mathcal P_{3m}$.
\end{itemize}
The force and the torque follow from~\eqref{eq:force-convention} and~\eqref{eq:torque}: $\Fsp=M^\star\nabla T|_{\vr^\star}$. This is the form~\eqref{eq:coplanar-l2} of Chapter~\ref{ch:coplanar} written with $|\kb|\cos(\cdot-\epsilon)=\Re(\kb\,e^{\I\,\cdot})$; it equals~\eqref{eq:impl-single} at $i=0$ to rounding error.

\paragraph{When to use it.} It costs $2(2q_{\max}+1)$ terms instead of $9(2q_{\max}+1)$. Its error with respect to~\eqref{eq:impl-single} is linear in the obliquity. For any Love numbers with $|\kb_2|\le\kappa_{\max}$ (the assumption of the truncation guarantee~\eqref{eq:guarantee}),
\begin{equation}
  \|\vF_{\rm coplanar}-\vF_{\rm single}\|\le C(e)\;i\;\kappa_{\max}F_0(r^\star), \label{eq:coplanar-switch}
\end{equation}
where $C(e)$ is the worst case over the orbit, reached at apocentre: $C=2.0$ ($e=0$), $2.7$ ($e=0.05$), $3.6$ ($0.1$), $7.1$ ($0.2$), $13$ ($0.3$), $24$ ($0.4$), $55$ ($0.5$). Realistic rheologies stay well below (Maxwell, constant time lag: $C\le2.5$ at $e=0.49$). The coplanar form is used for $i\le i_s=10^{-7}$\,rad: the jump at the switch is then below $5.5\times10^{-6}\,\kappa_{\max}F_0$ for all $e\le0.5$, i.e.\ below the truncation tolerance $10^{-5}$. It applies to a prograde orbit only ($\hat{\mathbf h}\cdot\hat{\mathbf s}>0$); near $i=\pi$ the surviving terms are $(0,1)$ and $(2,2)$, and the general form~\eqref{eq:impl-single} is used.

\section{Dissipated power}\label{sec:power}
$N$-body codes often output the power dissipated by the tide in the planet. For the constant-time-lag model it is $-(\mathbf F_{\rm diss}\cdot\mathbf v+\boldsymbol\Gamma\cdot\mathbf s)$, with the dissipative part of the force only \citep{Bolmont2015}. The same definition is used here: the dissipative part is the part of the force carried by $\Im\kb_2$,
\begin{equation}
  \boxed{P_{\rm diss}=-\big(\mathbf F_{\Im}\cdot\mathbf v+\boldsymbol\Gamma_{\Im}\cdot\mathbf s\big),} \label{eq:power}
\end{equation}
where $\mathbf v=\mathbf v_{\mathrm p}-\mathbf v_\star$, and $\mathbf F_{\Im}$, $\boldsymbol\Gamma_{\Im}$ are the force~\eqref{eq:impl-force} and the torque~\eqref{eq:torque} computed with every $\kb_2(\bar\sigma_{2mpq})$ replaced by $\I\,\Im\kb_2(\bar\sigma_{2mpq})$. Since the force is linear in $(\Re\kb_2,\Im\kb_2)$, $\mathbf F_{\Im}$ and $\boldsymbol\Gamma_{\Im}$ are obtained by a second accumulation of the same terms.
\begin{itemize}
  \item \textbf{Orbit average.} Over an orbit with constant rates,
  \begin{equation}
    \langle P_{\rm diss}\rangle=\frac{\G M^{\star2}R^5}{a^6}\sum_{m,p,q}c_m\,F_{2mp}^2(i)\,G_{2pq}^2(e)\;\bar\sigma_{2mpq}\,\big(-\Im\kb_2(\bar\sigma_{2mpq})\big)\ \ge0, \label{eq:power-mean}
  \end{equation}
  non-negative since $\sigma\,\Im\kb_2(\sigma)\le0$ for a lag.
  \item \textbf{Instantaneous value.} $P_{\rm diss}$ is not positive at every instant for a general rheology (minimum $-0.8\,\%$ of the mean in an Andrade case at $e=0.3$, against $-860\,\%$ for the power of the total force). The split into $\Re$ and $\Im$ parts is a convention, the one that extends the constant-time-lag output: in that limit~\eqref{eq:power} equals the constant-time-lag power up to $\mathcal O\big((n\Delta t)^2\big)$. For an elastic planet it is exactly zero.
\end{itemize}

\section{Algorithm}\label{sec:algorithm}
For each planet with the tide active, at each force evaluation:
\begin{enumerate}
  \item \textbf{Geometry} (Section~\ref{sec:geometry}): $\hat{\mathbf s}$, $\Op$, $\hat{\mathbf h}$; nodal frame $\hat{\mathbf N},\hat{\mathbf Y}$; $i$, $\hat{\mathbf W}$, $\hat{\mathbf N}'$; $a$, $e$, $\bar n_{\rm osc}$, $u$, $\omega$, $f$, $\M$, $u_\M$; $r^\star$, $\rho$, $c^\star$.
  \item \textbf{Mean rates} (Section~\ref{sec:rates}): $\bar n$, $\dot{\bar\lambda}$, $\dot{\bar\Omega}$ from the history (read only).
  \item \textbf{Kaula functions}: $q_{\max}(e)$ from the table of Section~\ref{sec:truncation}; $F_{2mp}(i)$ for $m,p=0,1,2$; $G_{2pq}(e)$ for $p=0,1$ and $|q|\le q_{\max}$, from~\eqref{eq:G-table} for $e\le0.5$.
  \item \textbf{Terms}: for each term $(m,p,q)$ with $F_{2mp}(i)\neq0$, compute $\bar\sigma_{2mpq}$~\eqref{eq:sigma-regular}, $\kb_2(\bar\sigma)$~\eqref{eq:love-impl}, $\mathcal K_{2mpq}$~\eqref{eq:K2}, $\mathcal E_{mpq}$~\eqref{eq:Ecal}; accumulate the single form~\eqref{eq:impl-single}, or the coplanar form~\eqref{eq:coplanar-nbody} if $i\le i_s$.
  \item \textbf{Output}: $\Fsp$~\eqref{eq:impl-force} and $\boldsymbol\Gamma$~\eqref{eq:torque}, and their dissipative parts for~\eqref{eq:power}.
\end{enumerate}
At each \emph{accepted} step, update $\lambda$ and $\Omega$, recompute the window~\eqref{eq:window}, and append $(t,\lambda,\Omega,n_{\rm osc})$ to the history.

\paragraph{Exact savings.} Two identities halve the cost of the Love numbers and of the phases, without changing the result beyond rounding:
\begin{itemize}
  \item for $m=0$, $\bar\sigma_{201,-q}=-\bar\sigma_{201q}$ and $\bar\sigma_{202q}=-\bar\sigma_{200,-q}$ exactly, and $\kb_2(-\sigma)=\kb_2(\sigma)^*$ (hypothesis H1): one evaluation of the spectrum per pair;
  \item $e^{\I q\M}$ is obtained by successive products of $e^{\I\M}$ (rounding error $\approx|q|\times10^{-16}$).
\end{itemize}

\paragraph{Cost.} At most $9\times(2q_{\max}+1)$ terms: 45 for $e\le0.001$, 63 for $e\le0.02$ and 765 for $e=0.5$. In the coplanar form, $2\times(2q_{\max}+1)$. Each term needs one interpolation in the spectra and one complex product.

\section{Validation tests}\label{sec:validation}
Any implementation should pass the following tests. Each one checks a property of the exact physics, independent of the Kaula derivation. The reference implementation (Python) and Posidonius pass all of them (Section~\ref{sec:pos-tests}).
\begin{description}
  \item[V1 Circular equatorial orbit, constant time lag.] The closed forms of Section~\ref{sec:physical-check}:
  \[
    F^{\mathrm p\to\star}_r=-\frac{3k_2\G M^{\star2}R^5}{a^7}\Big[\tfrac14+\tfrac34\cos2(\Op-n)\Delta t\Big],\qquad
    F^{\mathrm p\to\star}_\varphi=\frac{3k_2\G M^{\star2}R^5}{2a^7}\sin2(\Op-n)\Delta t,
  \]
  and $F_\theta=0$.
  Agreement $5\times10^{-16}$.
  \item[V2 Constant time lag, eccentric and oblique.] With $\kb_2=k_2e^{-\I\sigma\Delta t}$ and Kepler rates, the force equals the exact force of the delayed bulge: the bulge points to the star position $\Delta t$ earlier, rotated with the planet by $\Op\Delta t$. Agreement $10^{-10}$ with $|q|\le40$.
  \item[V3 Small time lag.] The difference with the Hut force \citep{Hut1981} decreases as $\Delta t$.
  \item[V4 Elastic planet.] With a constant real $k_2$, for any rates and any obliquity, $\boldsymbol\Gamma=0$ and
  \[
    \Fsp=-3k_2\G M^{\star2}R^5(r^\star)^{-7}\,\hat{\vr}.
  \]
  \item[V5 Rotation of the inertial axes.] Rotating $(\vr,\mathbf v,\mathbf s)$ rotates $\Fsp$ and $\boldsymbol\Gamma$ in the same way. This test detects any hidden reference to the inertial axes or to an inertial plane.
  \item[V6 Change of length unit.] Multiply all lengths by $\lambda$ and $\G$ by $\lambda^3$, which keeps the time unit. Then $\Fsp\to\lambda\Fsp$ and $\boldsymbol\Gamma\to\lambda^2\boldsymbol\Gamma$. This test detects dimensional errors.
  \item[V7 Obliquity $\to0$.] The force is continuous at $i=0$, with $|\vF(i)-\vF(0)|=\mathcal O(i)$, and independent of the node direction chosen at $i=0$.
  \item[V8 Signs.] For a coplanar orbit and a constant time lag, the orbit-averaged $\boldsymbol\Gamma\cdot\hat{\mathbf s}$ is negative for $\Op>n$ and positive for $\Op<n$. For an oblique spin, the averaged torque equals the averaged Hut torque to $\mathcal O(\Delta t)$.
  \item[V9 Angular momentum.] In a two-body integration with the instantaneous torque~\eqref{eq:torque}, $\mathbf L_{\rm orb}+I\mathbf s$ is conserved up to the integration error of the scheme.
\end{description}

\section{Assumptions and limits}\label{sec:limits}
\begin{itemize}
  \item \textbf{Physics:} hypotheses H1--H3, response at the mean frequency. The result is exact for a constant time lag. For a general rheology, the error is of order $\partial_\sigma\kb_2\cdot\delta\sigma$, where $\delta\sigma$ is the neglected fast modulation of the rates.
  \item \textbf{Elements:} two-body osculating elements with $\mu_{\rm orb}=\G(M^\star+M_{\mathrm p})$, referred to the instantaneous equator.
  \item \textbf{Degree:} degree two only. The neglected degree-three terms are of relative order $(R/a)^2$.
  \item \textbf{Scope:} the tide in the planet. The tide in the star uses the same equations with the roles of the bodies exchanged (Chapter~\ref{ch:star}).
  \item \textbf{Eccentricity:} $e\le0.5$ with the truncation table above. Beyond, the code falls back to the largest $q_{\max}$ of the table ($42$) and Tisserand's formula, without guarantee.
  \item \textbf{Spin--orbit synchronism.} Hypothesis H3 fails for a term whose mean frequency is zero. At the 1:1 spin--orbit resonance, $\bar\sigma_{2200}=2(\dot{\bar\lambda}-\Op)$ uses the instantaneous spin rate $\Op$, while the spin librates with the orbital period: $\bar\sigma_{2200}$ changes sign twice per orbit. Its Love number is read at a frequency that oscillates around zero: $\Im\kb_2$ changes sign with $\bar\sigma_{2200}$, and jumps if the spectrum does not tend to zero at low frequency. The torque of this term then behaves like a dry friction (stick--slip): the spin equation becomes stiff and discontinuous, and the secular rates obtained in a run depend on the integration step. With an obliquity, the terms $(1,0,-1)$, $(1,1,1)$ and $(2,1,2)$ also have frequencies close to zero. A treatment of the response in the time domain (the bulge as a state variable with its own relaxation equations) would be needed in this regime; it is not part of this document.
\end{itemize}
